\documentclass[a4paper]{article}

%% Language and font encodings
\usepackage[english]{babel}
\usepackage[utf8x]{inputenc}
\usepackage[T1]{fontenc}
\usepackage{multirow}
%% Sets page size and margins
\usepackage[a4paper,top=3cm,bottom=2cm,left=3cm,right=3cm,marginparwidth=1.75cm]{geometry}

%% Useful packages
\usepackage{amsmath}
\usepackage{graphicx}
%%\usepackage[colorinlistoftodos]{todonotes}
\usepackage[colorlinks=true, allcolors=blue]{hyperref}


\newcommand*{\bord}{\multicolumn{1}{c|}{}} 
\newcommand*{\nobord}{\multicolumn{1}{c}{}}

\title{Bootstrapping charges via Supersymmetry}
\author{Alejandro Rivero\footnote{arivero@unizar.es}}

\begin{document}
\maketitle

\begin{abstract}
A self-referencing postulate in the Chan-Paton
factors of a string allows to fix some of the freedom in the choosing of
standard-model-like vacua.

\hrulefill
\hrulefill
\end{abstract}

With the failure to find scalar superpartners of the currently known
quarks and leptons, it becomes of interest to explore other 
representations of supersymmetry (eg \cite{Jourjine:2012hv,Armori}). 
In this letter we find if not a representation of susy, at least
an interesting representation of three generations of scalars that 
could be used to form SM supermultiplets. 

\subsubsection*{The view from above}
Consider $SO(32)$. We can factor it as flavour times colour, $G^F \times SU(3)$, following the
procedure of  \cite{GellMann:1976pg} where it is listed as case 4, equation
2.18:
\begin{eqnarray*}
(32 \times 32)_A= [1,1,1^c] + {\bf (1,24,1^c) }+ (1,1,1^c) 
        +(1,1,8^c) + (2,5,3^c) + (2,\bar 5, \bar 3^c) \\
      +{\bf [1,15,\bar3^c]} + {\bf [1, \bar {15}, 3^c]} 
      +[1,10,6^c] + [1, \bar 10, \bar 6^c] 
\end{eqnarray*}

where the derived flavour group is $G^F=SO(2) \times U(5) \times U(1)$.  In
fact the same structure plus an extra $(1,28,8^c)$ can be derived branching down
via 
$$SO(32) \supset SO(30) \times U(1) \supset SU(15) \times U(1) \times U(1)
\supset SU(5) \times SU(3) \times U(1) \times U(1) $$ 
from where, following the branching rules in \cite{Yamatsu:2015npn} we get
 \begin{eqnarray*}
{\bf 496 }=    (24, 8^c)_{0,0} 
   ⊕ (1,1^c)_{0,0} \oplus {\bf (24, 1^c)_{0,0}} \oplus (1,1^c)_{0,0} \\
  ⊕ (1, 8^c)_{0,0} ⊕   (5, 3^c)_{2,2} ⊕ (5, 3^c)_{2,−2} \oplus  (5, 3^c)_{-2,2} ⊕ (5, 3^c)_{-2,−2}\\ \oplus {\bf
   (15, 3^c)_{0,4} ⊕  ⊕ (15, 3^c)_{0,−4}} 
    ⊕ (10, 6^c)_{0,4} \oplus (10, 6^c)_{0,−4} 
\end{eqnarray*}
In this case a second $U(1)$ from the flavour part is explicitly raided out. Still another $U(1)$ charge can be got
by further breaking to $SU(5) \supset SU(3)\times SU(2) \times U(1)$.
Then the  $\bf 24$ descend as
$$
24 = (1, 1)(0) + (1, 3)(0) + (3, 2)(5) + (\bar 3, 2)(−5) + (8, 1)(0)
$$
and with $Q=\frac 15 Y$ it looks as three generations of scalar leptons. Furthermore,
the $\bf 15$ descends as
$$
15 = (1, 3)(−6) + (3, 2)(−1) + (6, 1)(4) \\
$$ and then extending the assignment of the $U(1)$ charges
the total $\bf 15+\bar {15}+24$ can be interpreted as three generations of scalar
quarks and leptons: 
$$
\begin{array}{l|c |r|r|c}  % a/b = 1/6 or = 2/3????
 irrep  & N& Y_1 & Y_2& Q= \frac 1{30}Y_1 - \frac15 Y_2 \\ % Q= -\frac 2{15}Y_1 - \frac15 Y_2
\hline
(6,1) & 6 & 4 &  4 &  -2/3  \\    
(3,2) & 6 & 4 & -1 &  +1/3  \\        
(1,3) & 3 & 4 & -6 & +4/3 \\           
(\bar 6,1) & 6 &-4 & -4 & +2/3\\         
(\bar 3,2) & 6 & -4 & 1 & -1/3\\          
(1,\bar 3) & 3 & -4 & 6 & -4/3\\                                                                                
(\bar 3,2)& 6 & 0 & -5 &     +1      \\   
(3,2) &6 &0 & 5 &   -1   \\
(1,1) & &&& \\
(1,3)&12 &0& 0& 0   \\
(8,1)& &&& \\
\end{array}        
$$      

In this sense, this is possibly the most straightforward way to obtain
three generations of standard model colour and electric charge out
of a string motivated group. This is probably left unnoticed because
one wants to get also the chiral charge, and then one looks for spinors 
in complex representations by exploring groups $SO(4n +2)$, which
excludes $SO(32)$.   

Isolated, 15+15+24 can be considered as a 54 of SO(10) or a 55 in Sp(10)
or O(10). Such representations are occasionally considered to organize "hidden mesons", for example recently in \cite{Arkani-Hamed:2016kpz,Yamatsu:2015npn} but 
the charge assignment in such cases is aimed to obtain colour and weak
charges as usual. 

In our case the table is interpreted as three generations of scalars with the correct electric
charge; plus three "half-generations" of an extra quark\footnote{it could be interesting to set also
the Q of this extra (1,3) to zero, but in order to do this we would need to grant some baryon number to (6,1) and (3,2)}.  To propose an identification for the components of
the $10$ of $SO(10)$ we can also branch the 10-plet of the SO(10) group first to $SU(5)$,
with $10 = 5 (2) + \bar 5 (-2)$ and then to $SU(3) \times SU(2)$,
via $5 = (3, 1)(2) + (1,2) (-3)$. We see that the 10-plet has three elements with $Q=\frac 2{30} - \frac 2{5} = -\frac 13$, two elements with $Q=\frac 2{30} + \frac 3{5} = \frac 23$  and then the corresponding opposite charges. Abusing a bit
on traditional notation, we can call $d,s,b$ to the
elements of the triplet and $u,c$ to the doublet, and draw the whole construction as [in figure]. It can be said that the bosons are actually being 
generated from pairing "quarks" or "antiquarks".

\subsubsection*{The view from below}

The goal of this paper is to point out a possible principle,
in this frame of open superstrings, to bootstrap this particular group
and charge from a self-referential argument. 

Lets consider here the $54=\bf 15+\bar {15}+24$ as produced by an unoriented bosonic string with $SO(10)$ or $Sp(10)$ Chan-Paton labels, and in a 
first step we delegate the
function  of colour to the string itself. We need unorientability of
the string in order to produce 
all the the 54 states forming a set of mesons, quarks and diquarks. 
Of course we could then again discard the unoriented states to reduce 
to a theory of mesons with SU(5) group, which in this case would be the three generations of scalar leptons.

Looking for a postulate for this structure, the idea that we propose
is to assume that there exists a
supersymmetric theory containing this same set of scalar states, and that there is 
a nuclear democracy in the labeling: the Chan-Paton charges of this theory are 
a subset of the quarks.  We can call them the light quarks or, indistinctly,
the preons of the theory.

So, consider $r$ quarks of $+2/3$ charge and $s$ quarks of charge $-1/3$. Our postulate is that they combine pairwise -at the extremes of an open string- to form consistent generations of squarks and sleptons. 

We can formulate this requisite in two ways:

- we can ask that the combination must build the same number of "up type" and "down type" diquarks. This is, that $$rs = s(s+1)/2 = 2n$$ and so $ s= 2 r -1$. Plus, we can ask $r$ to be even, to make sure we build a pair number of scalars of each charge (we want to be able to promote them to a supersymmetric theory in the future).

- or, we can go for a more general condition, that we will see implies the former one: we can ask that the total number of combinations must be an integer multiple of a single set, this is a multiple of $rs$.  In this case we look for integer positive solutions of:

$$rs + {s^2+s \over 2} + {r^2 + r \over 2} + (r+s)^2 -1 = K rs$$

and then $K=9$. This in turns implies that one generation is composed of $(K-1)/2$
Dirac-like tuples -including the neutrino-,
plus one extra state, with its antiparticle. This extra state can be
either neutral or charged with +4/3. 

In any case, either requirement fixes $s=2 r -1$ and then the allowed groups are 
$SO(2s+2r)=SO(6r -2)$, with $r$  even. The smallest possible group is,
as expected, SO(10), but bigger groups are possible. 

We could at this level to discuss if the representation is the
54 of SO(10) or the 55 of Sp(10) or just use O(10); either of them does
the job and in principle Sp(10) looks less ad-hoc. There is the inconvenience 
of an extra singlet that separates when branching down to SU(5). As this
branching is the same in all the cases, perhaps the best answer is
that at this level we are just using representations of $SU(r+s)$ and
that its collection in a $SO$ is just a way to keep the tabs.

Whan we need now is an extra postulate to fix $n_g=3$, if we are not
happy simply with being the smallest of the possible answers. 
To look for it, lets examine a
table of solutions, with the corresponding total of "standard model
generations" of scalars and its extra content, that we see adds always to $n_g$.
so one extra scalar in each generation.

$$
\begin{array}{cc|c|c|c|}
r & s & n_g & {r(r+1)\over 2 } & {r^2+s^2 -1 \over 2} - 2 n_g \\ 
\hline
2 & 3 & 3   &  3  & 0 \\
4 & 7 & 14  & 10  & 4 \\
6 & 11 & 33 & 21  & 12 \\
8 & 15 & 60 & 36 & 24 \\
%10 & 19 & 95 & 55 & 40 \\
%... \\
%16 & 31 & 248 & 136 & 112 \\
... \\
\end{array}
$$

From the table, it seems that the simplest extra postulate is to ask
for absence of excess neutral bosons. An equivalent requisite is to
ask that the number of neutral bosons to be equal to the number of 
charged ones, given that the later is always $2 r s$ as in the "quark" 
sector. 
 
 With this extra requisite, the solution is unique:
 $$ r=2, s=3, n_g=3, G=SO(10)$$
 

Lets now address colour. We already know that it can be recovered if one
considers to assign a SU(3) singlet to the 24 of SU(5) and SU(3) triplets and anti-triplets
to the $15$ and $\bar 15$. The decomposition should then be promoted to
something under $SO(10) \otimes SU(3)$ such as
$$ \bf ? \longrightarrow (15,3) + (\bar 15, \bar 3) + (24, 1) $$ 
What we need is some argument to pinpoint exactly $SU(3)$,
either as a preferred labeling or extracting it out of a $SU(3 + N)$ in large $N$.
%it could be interesting to know how to extract simply an U(1) for leptonic/barionic number 

We can argue, again self-referencially, that $N=3$ is derived from the fact that 
$$ \bf 3 \times 3 = 6 + 3 $$
so in some sense a composite $3$ is composed of itself. The mesonic
strings, on the other hand, go to the singlet from $$ \bf 3 \times \bar 3 = 8 + 1 $$ and the colouring paints a matrix

$${\bf 
\begin{array}{|cccccccccc|c} 
\hline
  &  &    & & \bord   &&&&& \\ \cline{1-1}
 \bord &&  && \bord&&&&& \\ \cline{2-2}
\nobord & \bord &\multicolumn{2}{r}{\bf \bar {3}}&\bord&&\multicolumn{3}{c}{\bf 1} & \\ \cline{3-3} 
\nobord&& \bord && \bord&&&&& \\ \cline{4-4}
\nobord &&& \bord &\bord&&&&& \\ \cline{5-10}
\nobord&&&& \bord &&&&& \\ \cline{6-6}
\nobord&&&&& \bord&&&& \\ \cline{7-7} 
\nobord&&&&&& \bord &\multicolumn{2}{c}{\bf 3}& \\ \cline{8-8} 
\nobord&&&&&&& \bord && \\ \cline{9-9}
\nobord&&&&&&&& \bord & \\    \cline{10-10}  
\end{array}} 
$$

where it can be seen more the process 
$$ oriented \to unoriented \to reoriented $$
that allows to extract out a set of mesons with Chan-Paton flavour charges
under $SU(r+s)$, the upper right box. 
This technique is used by Armori et al under the market 
name "planar orientifold" for the calculation of meson masses in a large $N$
theory generating an "Hadronic Susy". Note that the diquark, 
unoriented part, can not survive in the large $N$, making already a sort of confinement.

\subsubsection*{From bottom to the top}

Now we notice that colouring the scheme has expanded the solution 10-plet 
$$ dsb \; uc \  \bar d\bar s\bar b \; \bar u\bar c $$ to
a 30-plet $d_{rgb}s_{rgb}...$, and so from $SO(10)$ to $SO(30)$,
and we are really only one "uncolored preon" out of reaching $SO(32)$.

Even if we had no results about consistency of gauge theories in
strings, we could have considered, by symmetry reasons, to leave
the colour boxes to overflow the diagonal by adding singlets, this is, 
to consider $U(3)$ instead of $SU(3)$ and the $\bf 54+1$ of $SO(10)$ so that the
full composition, colour+flavour, of $U(3) \times O(10)$ should be something
adding as
\begin{eqnarray*}
\bf 
(15,3)+(15,6)+(\bar {15},\bar 3)+(\bar {15},6)+(24,8)+(24,1)+(1,8)+(1,1) \\
=(15,9)+(\bar {15}, 9)+(24,9)+(1,9)  
\\ =(54+1,9) \\ =495
\end{eqnarray*}

So we see that our super-bootstrap of a three family structure
when supplemented with three colours has about the right size
to potentially reach the Green-Schwarz cancellation.

%Permite el coloreado tambien decidir si es SO(x), Sp(x) o ninguno?

\subsubsection*{The chiral basis}(TO BE DONE, usease, esto es a lo
que le estoy dando vueltas ahora)


A clear inconvenient of our process is that it works in a non-chiral schema,
using electric charge and colour. We woild like to obtain similar results
for the chiral charges of the standard model, namely weak isospin
and weak hypercharge.

A first step could be to group our mesons in pairs associated to
 chiral supermultiplets. The simplest idea is to use a definition
 of left and rigt components as the sum and difference of particle/antiparticle.

$$d^L=d+\bar d$$
$$d^R=d-\bar d $$
$$...$$

(Here in principle we have a freedom to choose which combination
corresponds to each antiparticle but we assume that any choosing can be
absorbed by the flavour symmetry)

Correspondingly, we would construct "complex scalars" from
combinations of scalar and antiscalar. A first detail of this idea is that it will mark some peculiarity for neutrinos,
as in our construction the composite scalar neutrals in the diagonal of the $\bf 24$
are they own antiparticles, and then, if unmixed, they will fail to have a 
a "Right component". This could be a feature.

... ... 

Ideally, we should be able to fold our $54$ into a left-right
sum of $27 \oplus 27$ and then add the coloring, as we did
in the previous case, the expectation being again to build some
$(9,27)^L + (9,27)^R  = 243 + 243$ that allows us to guess
a valid decompositio. Lacking that, we could try to force our way
from above, possibly $E8 \times E8$ to $(SU(3) \times E6)
\times (SU(3) \times E6)$. A main issue is that we need to
understand why colour should act separately in both parts, and
we need to have triplets and singlets of colour in the 
adequate representations, as it happened in the $SO(32)$ descent.


For reference (working draft, again) lets list all the composites:

\begin{itemize}
\item In the 24 +1:
\begin{itemize}
\item charge +1, irrep $(3,2)$ 
$$\begin{array}{ccc} 
                 u \bar d & u\bar s & u \bar b \\ 
                 c \bar d & c\bar s & c \bar b \\
                \end{array}$$
\item charge -1, irrep $\bar (3,2)$
 $$\begin{array}{ccc} 
                 \bar u  d &\bar u s &\bar u  b \\ 
                 \bar c  d &\bar c s &\bar c  b \\
                \end{array}$$
\item charge 0, sum $ (1,3) + (8,1) + (1,1) + 1 $
  $$\begin{array}{ccccc}
        u \bar u & c \bar c & u \bar c & c \bar u \\
        d\bar s & d\bar b & s\bar b & d\bar d \\
        s\bar d & b\bar d & b\bar s & s\bar s & b\bar b \\
        \end{array}$$
\end{itemize}
\item From the $\bar {15} \oplus 15$ :
\begin{itemize}
\item charge +4/3 and -4/3
$$\begin{array}{ccc} 
                 u u &  c c &  c u \\ 
                 \bar  u\bar u &  \bar  c\bar c &  \bar c\bar u \\
                \end{array}$$
\item charge +2/3
$$\begin{array}{ccc} 
                 \bar d \bar d & \bar s\bar s & \bar b \bar b \\ 
                 \bar s \bar d & \bar b\bar s & \bar d \bar b \\
                \end{array}$$
                
\item charge -2/3
$$\begin{array}{ccc} 
                 d d &  s s &  b b \\ 
                 s d &  b s &  d b \\
                \end{array}$$
\item charge -1/3
$$\begin{array}{ccc} 
                 \bar u \bar d & \bar u\bar s & \bar u \bar b \\ 
                 \bar c \bar d & \bar c\bar s & \bar c \bar b \\
                \end{array}$$
                \item charge +1/3
$$\begin{array}{ccc} 
                  u  d & u s & u b \\ 
                  c  d & c s & c b \\
                \end{array}$$
\end {itemize}
\end{itemize}

 
... ...\dots{}

\subsubsection*{Conclusion}


In conclusion, lets review what we have got. We have started from the
postulate that a bosonic unoriented string could be constructed such that:

\begin{itemize}
\item The "preons" labeling the string are of two kinds according electric
charge: "up" and "down".
\item In the diquark sector, the number of "up" diquarks is equal to the number
of "down" diquarks.
\item In the meson sector, the number of charged mesons is equal to the number
of neutral mesons
\end{itemize}
 
We call this arrangement "a supersymmetrical bootstrap" because
in a field theory of quarks and mesons 
the preons should be a subset of the quarks of the theory,
assuming that such quarks are recovered via susy from the bosonic
sector. 

The main conclusion is that there is a consistent assignment of charges,
and the only possible number of generations is three.

Then we add colouring compatible with the construction: a meson must be 
a singlet of colour, a diquark must be in the same representation
of colour than a (anti-)quark. This isolates $SU(3)$ and strongly signals
that the full group must be $SO(32)$

The complete program to build such model should imply to start from
this bosonic "hadron-diquark" string model, upgrade to superstrings
then going to D=10, and then analyze the branes corresponding to
Chan-Paton matrices. Note that it would still be true that 
removing coloured sectors the hadronic mesons of the standard
model are recovered, so as a minimum the approach meets the
goals of the original program of 1970!

It is interesting that the number of quarks actually involved in 
the construction is one less than the number of quarks obtained
in the process, and that the one odd out is of "up" type. So it could
be expected that the process of building the model allows, or 
even forces, a mass for the top quark separated from
the other five.

Bootstrapping, as it is known, refers to the possibility of determining
some parameters -usually of an scattering process- via some set of consistency 
conditions and a self-referential mechanism -usually a duality-. This was
expected to generate some non-linearity pinpointing special places of the
space of parameters. Back in 1970, such program was still considered unable to pinpoint the
quantum numbers of a model, as Chew himself remarked in an article
in Physics Today:
{\it The unfriendly question raised most
often ... is how self-consistency can possibly be expected to generate "internal quantum numbers" such as hypercharge and baryon number..., but how can you bootstrap a symmetry? A conceivable response is that symmetries (or the associated quantum numbers) are related to particle multiplicities, and the non-linear unitarity condition responds to the number of different particles.} 
Nowadays the progress on the analysis of anomalies makes us more
knowledgeable so that today it is more likely 
to fix the particle content of gauge group via consistency conditions. What we have
exemplified here is that a combination of attacks from below and from above
can allow us to select some particular vacuum out of the landscape of 
string, and then specific generations of standard model like families. 

\begin{thebibliography}{99}

%\cite{Arkani-Hamed:2016kpz}
\bibitem{Arkani-Hamed:2016kpz}
  N.~Arkani-Hamed, R.~T.~D'Agnolo, M.~Low and D.~Pinner,
  %``Unification and New Particles at the LHC,''
  JHEP {\bf 1611} (2016) 082
  doi:10.1007/JHEP11(2016)082
  [arXiv:1608.01675 [hep-ph]].
  %%CITATION = doi:10.1007/JHEP11(2016)082;%%
  %3 citations counted in INSPIRE as of 27 Jul 2017

%\cite{Jourjine:2012hv}
\bibitem{Jourjine:2012hv}
  A.~N.~Jourjine,
  %``Scalar supersymmetry,''
  Phys.\ Lett.\ B {\bf 727} (2013) 211
  doi:10.1016/j.physletb.2013.10.036
  [arXiv:1208.1867 [hep-ph]].
  %%CITATION = doi:10.1016/j.physletb.2013.10.036;%%
  %6 citations counted in INSPIRE as of 31 Jul 2017

%\cite{GellMann:1976pg}
\bibitem{GellMann:1976pg}
  M.~Gell-Mann, P.~Ramond and R.~Slansky,
  %``Color Embeddings, Charge Assignments, and Proton Stability in Unified Gauge Theories,''
  Rev.\ Mod.\ Phys.\  {\bf 50} (1978) 721.
  doi:10.1103/RevModPhys.50.721
  %%CITATION = doi:10.1103/RevModPhys.50.721;%%
  %338 citations counted in INSPIRE as of 27 Jul 2017

%\cite{Harigaya:2016eol}
\bibitem{Harigaya:2016eol}
  K.~Harigaya and Y.~Nomura,
  %``Hidden Pion Varieties in Composite Models for Diphoton Resonances,''
  Phys.\ Rev.\ D {\bf 94} (2016) no.7,  075004
  doi:10.1103/PhysRevD.94.075004
  [arXiv:1603.05774 [hep-ph]].
  %%CITATION = doi:10.1103/PhysRevD.94.075004;%%
  %21 citations counted in INSPIRE as of 27 Jul 2017
  
  %\cite{Yamatsu:2015npn}
\bibitem{Yamatsu:2015npn}
  N.~Yamatsu,
  %``Finite-Dimensional Lie Algebras and Their Representations for Unified Model Building,''
  arXiv:1511.08771 [hep-ph].
  %%CITATION = ARXIV:1511.08771;%%
  %7 citations counted in INSPIRE as of 27 Jul 2017
\end{thebibliography}
\end{document}





And each representation goes down to (flavour) $SU(3) \times SU(2) \times U_2(1)$

\begin{eqnarray*}
15 =& (1, 3)(−6) + (3, 2)(−1) + (6, 1)(4) \\
24 =& (1, 1)(0) + (1, 3)(0) + (3, 2)(5) + (\bar 3, 2)(−5) + (8, 1)(0)
\end{eqnarray*}


In this model we look only to electric charge. Setting the two $U(1)$ charges to $Q_1+Q_2=-\frac 16$, $Q_2=-1/5$ we can produce a total combination $Q$ as
in the following table. Note that $N$ here is the number of states. 

$$
\begin{array}{l|c |r|r|c}  % a/b = 1/6 or = 2/3????
 irrep  & N& Y_1 & Y_2& Q= \frac 1{30}Y_1 - \frac15 Y_2 \\ % Q= -\frac 2{15}Y_1 - \frac15 Y_2
\hline
(6,1) & 6 & 4 &  4 &  -2/3  \\    
(3,2) & 6 & 4 & -1 &  +1/3  \\        
(1,3) & 3 & 4 & -6 & +4/3 \\           
(\bar 6,1) & 6 &-4 & -4 & +2/3\\         
(\bar 3,2) & 6 & -4 & 1 & -1/3\\          
(1,\bar 3) & 3 & -4 & 6 & -4/3\\                                                                                
(\bar 3,2)& 6 & 0 & -5 &     +1      \\   
(3,2) &6 &0 & 5 &   -1   \\
(1,1) & &&& \\
(1,3)&12 &0& 0& 0   \\
(8,1)& &&& \\
\end{array}        
$$            


 
$$ {{\bf 54} \longrightarrow
\begin{array}{|cccccccccc|c} 
\hline
  &  &    & & \bord   &&&&& \\ \cline{1-1}
 \bord &&  && \bord&&&&& \\ \cline{2-2}
\nobord & \bord &\multicolumn{2}{r}{\bf \bar {15}}&\bord&&\multicolumn{3}{c}{\bf 24} & \\ \cline{3-3} 
\nobord&& \bord && \bord&&&&& \\ \cline{4-4}
\nobord &&& \bord &\bord&&&&& \\ \cline{5-10}
\nobord&&&& \bord &&&&& \\ \cline{6-6}
\nobord&&&&& \bord&&&& \\ \cline{7-7} 
\nobord&&&&&& \bord &\multicolumn{2}{c}{\bf 15}& \\ \cline{8-8} 
\nobord&&&&&&& \bord && \\ \cline{9-9}
\nobord&&&&&&&& \bord & \\    \cline{10-10}  
\end{array}} 
\longrightarrow
\begin{array}{c|cccccccccc|}
 &d&s&b&u&c&\bar d&\bar s&\bar b&\bar u&\bar c \\
\hline
d &  & \multirow{2}{*}{$-\frac 23$} &  \bord  & & \bord   &&&\bord && \\ \cline{2-2}
s & \bord &  & \bord & \multicolumn{2}{c|}{\frac 13} &&0&\bord&\multicolumn{2}{c|}{-1} \\ \cline{3-3}
b && \bord &\bord&&\bord&&&\bord&& \\ \cline{4-11} 
u &&& \bord && \multicolumn{1}{c|}{\multirow{2}{*}{$\frac {-4}3$}} &&\multirow{2}{*}{+1}&\bord&\multicolumn{2}{c|}{\multirow{2}{*}{0}} \\ \cline{5-5}
c &&&& \bord &\bord&&&\bord&\multicolumn{2}{c|}{} \\ \cline{6-11}
\bar d &&&&& \bord &&\multirow{2}{*}{$\frac 23$}&\bord&& \\ \cline{7-7}
\bar s &&&&&& \bord&&\bord& \multicolumn{2}{c|}{-\frac 13} \\ \cline{8-8} 
\bar b &&&&&&& \bord &\bord&& \\ \cline{9-11} 
\bar u &&&&&&&& \bord &&\multirow{2}{*}{$\frac 43$} \\ \cline{10-10}
\bar c &&&&&&&&& \bord & \\    \cline{11-11}  
\end{array} 
$$
%%posiblemente habria sido mas astuto poner una caja de cuatro y luego cada cuadrante por separado, ajustando anchos


 Note
that honestly we are still in the world of the bosonic string; we
can take this number as sheer coincidence or as signal that the
"sBoostrap" postulate actually has some susy in it. In any
case, it isolates a pathway for model building with
exactly three generations.

Actually the sum is a bit of luck; a more accurate decomposition,
for instance, is to branch out from $SO(32)$ to $SU(5)\times SU(3)$ via 
say $SO(30)$ and $SU(15)$. Such pathway decomposes:

\begin{eqnarray*}\bf
496 = (435)_0 ⊕ (30)_2 ⊕ (30)_{−2} ⊕ (1)_0 \\ \bf
=   (224)_{0,0} ⊕ (105)_{0,4} ⊕ (105)_{0,−4} ⊕ (1)_{0,0} \oplus 
   (15)_{2,2} ⊕ (15)_{2,−2} \oplus  (15)_{-2,2} ⊕ (15)_{-2,−2} \\ \bf
=  (24, 8)_{0,0} ⊕ (24, 1)_{0,0} ⊕ (1, 8)_{0,0} ⊕ \\ \bf
   (15, 3)_{0,4} ⊕ (10, 6)_{0,4} ⊕ (15, 3)_{0,−4} ⊕ (10, 6)_{0,−4} ⊕ (1)_{0,0} \oplus \\ \bf
   (5, 3)_{2,2} ⊕ (5, 3)_{2,−2} \oplus  (5, 3)_{-2,2} ⊕ (5, 3)_{-2,−2} \\ \bf
\end{eqnarray*}


In the current context, as we are not including the weak
force in the model, there is no value on locating
a concrete descent. Also, we ignore if there is some 
more direct way to relate $SO(10)$ to $SO(32) = SO(2^{D/2})$;
is is a bit intriguing that adding colour seems to be enough.



\end{document}

%hablando de 27:  54/2 == 27
  asi que igual podemos doblar la 54 sobre si misma? antes de aplicar el color?


%oriented ---> unoriented --> reoriented
 SU(2n)   ----> SO(2n) or Sp(n) ---> SU(n)

54 = (1, 1) + (5, 1) + (3, 7) + (1, 27)
7 = 1 + 6
27 = 1 + 6 + 20'


54 = 1 + 9 + 44
  9 = (1, 3) + (6, 1)
  44 = (1, 1) + (1, 5) + (6, 3) + (20', 1)

  9 = 1 + 8v
  44 = 1 + 8v + 35v

54 = (1, 1, 1) + (3, 3, 1) + (2, 2, 6) + (1, 1, 20'
)


(15,3) + (24,1) + (15,3) === 45 + 24 + 45 == 114
